\documentclass[12pt,letterpaper]{article}

% - package setup
\usepackage[doublespacing]{setspace}
\usepackage{csvsimple}

\usepackage{xcolor}
\definecolor{light-gray}{gray}{0.95}
\newcommand{\code}[1]{\colorbox{light-gray}{\texttt{#1}}}

% - title setup
\title{Midterm Project}
\author{Crystal Mandal}
\date{last edited \today}

\begin{document}

\maketitle

\section{Question 1:}
\subsection{Flat Random Sounds}
The music generated with the requisite parameters (by default in intervals of 40 sounds in
three voices each, at least ``16 measures" long) is displayed below:

\lilypondfile[staffsize=12, quote]{../lilypond/flat-random.ly}

\section{Question 2:}
\subsection{Markov Bank}
The bank of pitches (default the first 12 sounds) used in markov chain generation is:

\lilypondfile[staffsize=12, quote]{../lilypond/markov-bank.ly}

\subsection{Markov Chain Table}
The table generated for the Markov Chain is:

\csvautotabular{../../python/markov-table.csv}

\subsection{Markov Chain}
The music generated with the requisite parameters (by default in intervals of 20 sounds, at least "16 measures" long)
by Markov Chain is displayed below:

\lilypondfile[staffsize=12, quote]{../lilypond/markov-chain.ly}

\section{Methodology}
The music is generated by the \code{ midterm.py } script in the
\code{ ./python } folder. The length of the \textbf{Flat Random Sounds},
the number of voices, the length of the \textbf{Markov Bank}, and the
length of the \textbf{Markov Chain} are all configurable.

To generate a new \code{ midterm.pdf } document, run the script
\code{ ./generate.sh  }. The usage for this command is detailled
by running \code{ ./generate.sh -h  } or \code { ./generate.sh --help  }

The generated document is available as \code{assignment-print.pdf} in the
root directory.

\end{document}
